\documentclass[11pt,a4paper]{article}
\usepackage{fontspec}
\setmainfont{DejaVu Sans}[Scale=0.90]
\usepackage{amsmath,amssymb}
\usepackage[margin=2.2cm]{geometry}
\usepackage{booktabs,array,tabularx}
\usepackage{graphicx}
\graphicspath{{../figures/}}
\usepackage{enumitem}
\usepackage{titlesec}
\usepackage{parskip}
\titlelabel{\thetitle.\quad}
\titleformat{\section}{\normalfont\large\bfseries}{\thesection.}{0.6em}{}
\titleformat{\subsection}{\normalfont\normalsize\bfseries}{\thesubsection}{0.6em}{}
\titlespacing*{\section}{0pt}{15pt}{5pt}
\titlespacing*{\subsection}{0pt}{11pt}{4pt}
\setlist[itemize]{leftmargin=1.4em,itemsep=1pt,topsep=3pt}
% Preserve fonts/margins; permit a final line-breaking pass for long prose.
\setlength{\emergencystretch}{2em}
\begin{document}
\begin{center}
{\Large\bfseries Intégration résonante et conservation}\\[4pt]
{\normalsize Mesure de collision et représentation entropique finie}\\[2pt]
{\small Partie I --- 28 septembre 2026}\\[2pt]
{\small Résultats finis sous hypothèses --- quatre revues indépendantes terminées (2 octobre 2026)}
\end{center}

\section{Question et exclusions}
La note 22 sépare le comptage des amplitudes et les hypothèses cinétiques.
Il faut maintenant intégrer la conservation d'énergie sans transformer
un poids élargi en événement discrètement résonant.
Nous étudions une approximation cinétique scalaire déclarée, à température
strictement positive, avec des énergies positives et des surfaces
résonantes régulières. Les zéros acoustiques, croisements, cohérences
dégénérées et points critiques demandent une analyse distincte.
Aucun opérateur AlN convergé n'est construit ici.

Après élimination d'une impulsion modulo un vecteur réciproque, posons
$\Delta=\varepsilon_p-\varepsilon_a-\varepsilon_b$.
Pour une amplitude intégrable $f$ et une surface régulière,
\[
 \int f(q)\delta(\Delta(q))\,dq
 =\int_{\Delta=0}\frac{f(q)}{|\nabla\Delta(q)|}\,dS.
\]
En coordonnées de graphe, le facteur correspondant est
$1/|\partial_{q_j}\Delta|$. Trouver des racines ne suffit pas :
il faut couvrir toutes les composantes et intégrer la bonne mesure.
L'identité suppose $\Delta\in C^1$ avec $\nabla\Delta\ne0$ sur la
surface (valeur régulière), une valeur de surface bien définie pour $f$
(par exemple $f$ continue au voisinage de $\{\Delta=0\}$) et
$f/|\nabla\Delta|$ intégrable sur la surface ; $\delta$ s'entend comme
limite de régularisations.

\section{Pourquoi réajuster des poids hors résonance ne suffit pas}
Pour le flux réversible usuel
\[
 F=k[n_p(1+n_a)(1+n_b)-(1+n_p)n_an_b],\qquad k\ge0,
\]
les populations de Bose aux énergies originales donnent
$F=kR(e^{-\beta\Delta}-1)$ avec $R=(1+n_p)n_an_b>0$.
L'énergie produite par cet événement vaut donc
\[
 \dot E=-\Delta F
       =kR\Delta(1-e^{-\beta\Delta})\ge0.
\]
Elle est strictement positive si $k>0$ et $\Delta\ne0$.
Les défauts de signes opposés ne se compensent pas.
C'est une obstruction pour ce flux, ces énergies et cette incidence ;
ce n'est pas une impossibilité pour toutes les représentations.
La positivité d'une matrice linéarisée ne prouve pas non plus
la stationnarité non linéaire de Bose.

\section{Une représentation faible compatible}
Soit $V=\operatorname{span}\{\phi_i\}$ un espace réel fini contenant
l'énergie physique : $\varepsilon=\sum_i e_i\phi_i$.
Introduisons une variable entropique et la population reconstruite :
\[
 \xi_\alpha(q)=\sum_i\alpha_i\phi_i(q),\qquad
 n_\alpha(q)=\frac{1}{e^{\xi_\alpha(q)}-1}.
\]
On travaille dans le domaine ouvert où $\xi_\alpha>0$ aux points
d'évaluation. Une quadrature volumique positive $Q_v$, dont la matrice
d'évaluation est de rang colonne plein, définit
\[
 U_i=Q_v(\phi_i n_\alpha),\qquad
 M_{ij}=Q_v[n_\alpha(1+n_\alpha)\phi_i\phi_j].
\]
Alors $dU/d\alpha=-M$ et $M$ est définie positive.

Pour chaque point de quadrature réactionnelle, posons
$\ell_r=\phi(p_r)-\phi(a_r)-\phi(\lambda_{b,r})$, avec poids $\omega_r>0$.
Écrivons $A_r=n_p(1+n_a)(1+n_b)$,
$B_r=(1+n_p)n_an_b$ et leur moyenne logarithmique
\[
 \Lambda_r=\frac{A_r-B_r}{\log A_r-\log B_r}>0,
 \qquad \Lambda_r=A_r\ \hbox{si }A_r=B_r.
\]
Comme $\log B_r-\log A_r=\ell_r^{\mathsf T}\alpha$, on obtient
\[
 K(\alpha)=\sum_r\omega_r\Lambda_r \ell_r\ell_r^{\mathsf T},
 \qquad \dot U=K\alpha,\qquad M\dot\alpha=-K\alpha.
\]
C'est une dynamique de moments ; les coefficients ne sont pas
automatiquement des nombres de phonons sur un maillage.
Numériquement, $A_r-B_r$ et $\Lambda_r$ s'évaluent à partir de
$\log A_r$ et $\log B_r$ : séparément, $B_r$ peut sous-déborder alors que
$e^{\ell_r^{\mathsf T}\alpha}$ déborde, et leur produit naïf donne NaN
pour des variables entropiques très séparées.

Pour l'entropie et l'énergie de quadrature,
\[
 S_Q=Q_v[(1+n)\log(1+n)-n\log n],\qquad E_Q=e^{\mathsf T}U,
\]
on a
\[
 \dot S_Q=\alpha^{\mathsf T}K\alpha\ge0,\qquad
 e^{\mathsf T}Ke=\sum_r\omega_r\Lambda_r\Delta_r^2.
\]
Sur de vraies racines $\Delta_r=0$, $Ke=0$, l'énergie $E_Q$
est conservée et $\alpha=\beta e$ est stationnaire.
À cet équilibre, la forme linéarisée symétrique est
$M_*^{-1/2}K_*M_*^{-1/2}$, positive, avec vecteur nul
$M_*^{1/2}e$.
Des invariants supplémentaires peuvent agrandir le noyau.

Ces identités établissent une structure finie et une solution locale
sur le domaine déclaré. Elles ne prouvent ni invariance globale du
domaine, ni convergence des trajectoires vers une équation continue,
ni conservation exacte de l'intégrale physique si $Q_v$ est approchée.
Si l'on utilise $\varepsilon_h$ et ses racines, l'invariant est celui
de cette énergie approchée, sauf justification supplémentaire.

La perte du domaine est effective, et non seulement non exclue.
Avec $\Phi=I_3$, des poids unités, $e=(1,1,1)$, une ligne parent
$(2,-1,0)$ et deux filles $(\tfrac12,-1,1)$, la résonance est exacte.
Partant de $\alpha=(1{,}0005;\,2;\,3)$, la variable entropique du parent,
évaluée hors des nœuds de volume, atteint zéro en un temps fini
$T\le1{,}11\times10^{-7}$ ; l'intégration numérique donne
$T\simeq9{,}46\times10^{-8}$. Jusque-là, $E_Q$ est conservée à
$10^{-15}$ près, $S_Q$ croît et les occupations de volume restent
positives, tandis que l'occupation du parent diverge. L'énergie et
l'entropie de quadrature ne contrôlent donc pas la réalisabilité
hors des nœuds.

\section{Contre-exemples discriminants}
L'interpolation doit porter sur la variable choisie.
À $\beta=\log2$, les populations de Bose aux énergies $1,3$
sont $1,1/7$. Au milieu, d'énergie $2$, interpoler les populations
donne $4/7$ au lieu de $1/3$. Pour deux filles d'énergie $1$,
le flux réduit à l'équilibre est alors
$3(4/7)-1=5/7$, malgré une interpolation exacte de l'énergie.
Interpoler $\xi$ donne la bonne valeur de Bose.
Hors équilibre, la même interpolation peut faire décroître $S_Q$ à
énergie conservée : un exemple exact donne $\dot E_Q=0$ et
$\dot S_Q\simeq-3{,}59\times10^{-3}$.

Dans les exemples analytiques réguliers des branches A et C,
omettre le Jacobien de surface produit une erreur de mesure de
$20\%$ tout en conservant les invariants testés.
Près d'un point critique, un simple balayage des changements de signe
peut manquer deux racines régulières : une sortie vide vérifie alors
les invariants de façon vacue.

Enfin, pour le modèle unidimensionnel $\Delta(x)=vx$, un échantillonnage
gaussien uniforme dépend du rapport $r=\sigma/(|v|h)$ :
\[
 h\sum_{k\in\mathbb Z}\delta_\sigma(vh(k+\theta))
 =\frac{1}{|v|}
  \left[1+2\sum_{m\ge1}e^{-2\pi^2m^2r^2}\cos(2\pi m\theta)\right].
\]
À $r$ fini, une erreur dépendant du décalage persiste.
La convergence uniforme en décalage dans ce modèle exige $r\to\infty$.
Ce n'est pas une règle universelle pour toutes les grilles et surfaces.
À maillage fixé, réduire $\sigma$ peut donner zéro ou une divergence.
Une faible dérive d'énergie ne certifie donc pas la convergence de mesure.

L'ordre de cette dérive exige lui-même une hypothèse. Près de
l'équilibre de Bose, un élargissement gaussien donne
\[
 \dot E_\sigma=\int\rho(s)\,\delta_\sigma(s)\,kR(s)\,s(1-e^{-\beta s})\,ds,
\]
où $\rho$ est la densité image, par $\Delta$, de la mesure pondérée.
Si $\rho$ et $R$ sont bornées et continues en $s=0$, alors
$\dot E_\sigma=\beta\sigma^2kR(0)\rho(0)+o(\sigma^2)$.
Pour un pli quadratique $\Delta=x^2$, $\rho(s)=s^{-1/2}$ n'est pas
bornée et l'ordre devient $\sigma^{3/2}$, de constante
$\beta kR(0)\int_0^\infty u^{3/2}\varphi(u)\,du\simeq0{,}4300\,\beta kR(0)$,
où $\varphi$ est la densité normale réduite.
Un ordre observé ne classe donc la surface qu'avec une densité bornée.

\section{Portée et prochaine étape matérielle}
Les contrôles analytiques, symboliques et numériques portent sur
des modèles réguliers explicitement définis. Les divergences critiques
observées ne classent pas toutes les singularités physiques.
Les poids natifs phono3py sont des combinaisons de canaux ; ils ne
constituent pas directement une quadrature globale d'événements positifs.

Conservation, signe de l'entropie et stationnarité de Bose ne
contraignent pas les taux : multiplier tous les poids $\omega_r$ par une
constante, ou vider la liste des réactions, les préserve tous. Seule une
consistance de l'action sur des fonctions non invariantes, stable par
raffinement, validerait les temps de collision.

Les filles identiques forment, sur une surface résonante régulière,
un ensemble de mesure nulle. Avec une densité de surface localement
bornée, la masse d'un tube de largeur $\ell$ autour de la diagonale est
$O(\ell^d)$ en codimension $d$ ; une densité seulement intégrable, par
exemple $|z|^{-a}$ avec $0<a<d$, ne donne que $\ell^{d-a}$. Un nœud de
quadrature situé sur cette diagonale représente une région voisine :
lui attribuer le facteur fini d'oscillateurs répétés de la note~22
modifierait la quadrature continue. Les identités de comptage et la
présente mesure ne se raccordent donc pas par les seuls labels.

L'application à l'AlN exige une dispersion cohérente, une énumération
des surfaces, leurs poids et symétries, une représentation des populations,
une étude de convergence et un régime cinétique justifié.
Le cadre utilise des idées établies de Galerkin entropique et de
quadrature implicite ; aucune nouveauté générale n'est revendiquée.

\section{Sources et reproduction}
\begin{itemize}
\item Fugallo et al., Phys.\ Rev.\ B \textbf{88}, 045430 (2013).
\par doi:10.1103/PhysRevB.88.045430.
\item Saye, SIAM J.\ Sci.\ Comput.\ \textbf{37}, A993--A1019 (2015).
\par doi:10.1137/140966290.
\item Abdelmalik et van Brummelen, J.\ Stat.\ Phys.\ \textbf{164},
77--104 (2016).
\par doi:10.1007/s10955-016-1529-5.
\end{itemize}
Les sources classiques éclairent les ingrédients mais ne prouvent
pas la convergence de leur combinaison pour l'opérateur Bose proposé.
Scripts, échecs, limites et quatre revues sont archivés dans
\texttt{theory/\allowbreak aln/\allowbreak resonance\_measure} ; le NaN
relevé par les revues est corrigé par l'évaluation logarithmique
(quatorze contrôles reproductibles).
\end{document}
